\documentclass[../../../include/open-logic-section]{subfiles}

\begin{document}

\olfileid{his}{set}{mythology}

\olsection{ઇતિહાસ કરતાં દંતકથા વધુ?}

એક સદીથી વધુ સમય પછી આ ઘટનાઓ તરફ
પાછું જોતાં, આ ચુકાદાઓને તપાસ્યા વિના
સ્વીકારી ન લેવા જોઈએ. પરિણામો ચોક્કસ
નવાં, રોમાંચક અને આશ્ચર્યજનક હતાં.
પરંતુ એ ખરેખર કેટલાં ચોંકાવનારાં હતાં?
અને શું તેમણે ખરેખર બતાવ્યું કે
ભૌમિતિક સહજજ્ઞાન પર આધાર ન રાખવો જોઈએ?

ચોંકાવનારી અસરના પ્રશ્ને,
\citet{Gouvea2011} દર્શાવે છે કે
કૅન્ટૉરે ડેડેકિન્ડને લખેલું પ્રસિદ્ધ
વાક્ય “\emph{je le vois, mais je ne le crois pas}”
ઘણી વાર તેના સંદર્ભથી જુદું કરીને લેવાય છે.
તેનો વધુ સંદર્ભ આ છે
(\citeauthor{Gouvea2011}માંથી ઉદ્ધૃત):
\begin{quote}
વિષય પ્રત્યેના મારા ઉત્સાહને લીધે
હું તમારી દયા અને ધીરજ પાસેથી
આટલી બધી અપેક્ષાઓ રાખું છું,
તે માટે મને માફ કરશો. હાલમાં
તમને મોકલેલી મારી વાતો મને
પોતાને પણ એટલી અણધારી,
એટલી નવી લાગે છે કે, માનનીય
મિત્ર, તેમની સાચાશ વિશે તમારો
નિર્ણય મળે ત્યાં સુધી મને
મનની શાંતિ મળી શકતી નથી.
તમે મારી સાથે સંમત ન થાઓ
ત્યાં સુધી હું એટલું જ કહી
શકું છું: \emph{je le vois, mais je ne le crois pas.}
\end{quote}
કૅન્ટૉર જાણતા હતા કે તેમનું
પરિણામ “એટલું અણધાર્યું,
એટલું નવું” છે. પરંતુ તેમને
પોતાનું પરિણામ ક્યારેય
\emph{અવિશ્વસનીય} લાગ્યું હોય
તે શંકાસ્પદ છે. \citeauthor{Gouvea2011}
દર્શાવે છે તેમ, તેમણે રજૂ કરેલી
સાબિતી ડેડેકિન્ડ તપાસી આપે
એટલું જ તેઓ ઇચ્છતા હતા.

ભૌમિતિક સહજજ્ઞાનના પ્રશ્ને,
પીઆનોએ પોતાની અવકાશ ભરતી
વક્રરેખા કોઈ ચિત્ર વિના
પ્રકાશિત કરી. પરંતુ હિલ્બર્ટે
પોતાની વક્રરેખા પ્રકાશિત
કરી ત્યારે પોતાનો હેતુ સમજાવ્યો:
વાચકો “નીચેના ભૌમિતિક
સહજજ્ઞાનની મદદ લે” તો
પીઆનોનું પરિણામ સમજવાની
સ્પષ્ટ રીત તેમને મળશે.
પછી તેમણે
\olref[his][set][pathology]{sec}માં
આપ્યાં છે તેવાં જ
\emph{ચિત્રો}ની શ્રેણી આપી.

વધુ સામાન્ય રીતે, પ્રકાશિત
સાબિતીઓમાં ચિત્રોનો વપરાશ
થોડો ઓછો થયો હોય, છતાં
ગણિતશાસ્ત્રીઓ સાબિતી કરતી
વખતે ચિત્રોનો \emph{વારંવાર}
ઉપયોગ કરે છે એ હકીકત
ટાળી શકાતી નથી. (સરળ
રીતે કહીએ તો, સારા
ગણિતશાસ્ત્રીઓ જાણે છે કે
ભૌમિતિક સહજજ્ઞાન પર
ક્યારે આધાર રાખી શકાય.)

એટલે આ વાતની અતિશયોક્તિ
માની ન લેવી; ઓછામાં ઓછું,
તપાસ્યા વિના તેનો સ્વીકાર
ન કરવો. વધુ જાણવા
\citet{Giaquinto2007} વાંચી શકો.

\end{document}
